\section{AI-agent research and implementation program}
\label{sec:program}

The AI agent is not used to replace physical validation.  Its role is to maintain
traceability across literature, material cases, calculations, code, and failed
hypotheses, while automating the repetitive parts of comparison and model
construction.

\subsection{Closed-loop workflow}

For each candidate material and phonon mode, the agent executes the following
auditable loop.

\begin{enumerate}
  \item \textbf{Read and extract.}  Search primary literature and record the
  material, mode eigenvector, $\vq$, electronic method, interaction parameters,
  definition and normalization of $g$, frequency at which the vertex is evaluated,
  and reported uncertainty.
  \item \textbf{Classify the perturbation.}  Project the DFPT perturbation,
  compute the charge, internal, and nonlocal channel weights of
  \cref{eq:channelstrength}, and attach a human-readable physical interpretation.
  \item \textbf{Select a benchmark.}  Choose the least expensive many-body method
  able to represent the active channel: UEG/vertex Monte Carlo for density,
  DFT+DMFT for local correlated channels, GW perturbation theory for nonlocal
  quasiparticle response, or an explicit model when interaction derivatives are
  central.
  \item \textbf{Generate a hypothesis.}  Propose a low-dimensional form of
  $\mathcal K_{ab}(\vq,\omega)$ with symmetry, Ward/gauge, locality, and causality
  constraints encoded before fitting.
  \item \textbf{Design a discriminating calculation.}  Select the mode or
  momentum point that maximizes disagreement between the current hypothesis and a
  plausible alternative, instead of merely adding more similar data.
  \item \textbf{Implement and test.}  Produce a versioned correction module,
  unit tests for algebraic invariants, and numerical regression tests against
  held-out high-level vertices.
  \item \textbf{Update or reject.}  Store negative evidence and revise the claim
  ledger.  The agent may not promote an empirical relation to an exact principle.
\end{enumerate}

Each extracted number must retain source, equation/table/figure location,
normalization convention, and whether it was read directly or inferred.  This is
essential because electron--phonon matrix elements can differ by phonon
normalization, spin factors, mode phase, orbital gauge, and screening convention.

\subsection{Minimal benchmark ladder}

The proposed ladder is designed to distinguish physical channels, not simply to
cover many chemical compositions.

\begin{table}[htbp]
\centering
\caption{Minimal benchmark ladder for developing and falsifying the method.}
\label{tab:benchmarks}
\begin{tabularx}{\linewidth}{@{}p{2.7cm}p{3.1cm}XX@{}}
\toprule
System & Active contrast & Reference level & Question answered \\
\midrule
Uniform electron gas & $q=0$ anchor versus $0<q\le2k_F$ density response &
Vertex diagrammatic Monte Carlo & Is the finite-$q$ density prior accurate and
what controls its residual? \\
SrVO$_3$ & Jahn--Teller versus breathing modes & DFT+DMFT & Does traceless orbital
weight predict the large mode-selective correction? \\
Doped CaCuO$_2$ & Half- versus full-breathing, static versus dynamic vertex &
DFT+DMFT & Can the dynamic-risk indicator detect failure of a static correction? \\
CoO & Correlated insulator and polar coupling & DFPT+$U$ / experiment & Does the
framework detect an invalid reference state and occupation-response channel? \\
Ba$_{1-x}$K$_x$BiO$_3$ or a simpler GWPT solid & Nonlocal quasiparticle response &
GW perturbation theory & Can a sparse nonlocal kernel reproduce the high-level
correction? \\
\bottomrule
\end{tabularx}
\end{table}

The UEG is the first diagnostic object because it isolates density physics.  It
must not be the only validation object, because success there cannot test internal
crystal channels.

\subsection{Accuracy, speed, and transparency metrics}

The target should be evaluated against three independent axes.

\paragraph{Accuracy.}
Report the error of complex matrix elements where phases are comparable, squared
couplings averaged on the relevant energy shell, and final observables.  A
practical initial threshold is to reduce the high-level-versus-DFPT error by at
least a factor of two while keeping the median corrected $|g|$ error below
10--15\% on held-out modes.  This threshold is a project criterion and may be
revised after the evidence matrix is complete.

\paragraph{Cost.}
Count the number of high-level displaced/response calculations, total core or GPU
hours, memory, and wall time relative to a dense many-body calculation of every
mode and $\vq$.  Speedup is meaningful only at comparable numerical convergence.

\paragraph{Transparency.}
For every prediction, output the operator decomposition, selected correction
level, training-source distance, uncertainty, and any violated invariant.  The
framework should expose why a mode was corrected rather than only returning a
number.

\subsection{Implementation layers}

The initial software can be separated into small testable modules:

\begin{enumerate}
  \item a reader for DFPT perturbation matrices and wavefunctions;
  \item localized-subspace construction and gauge checks;
  \item operator decomposition and Fermi-window channel weights;
  \item adapters for high-level reference data from UEG, DFT+DMFT, DFPT+$U$, and
  GWPT;
  \item constrained kernel fitting/interpolation with uncertainty;
  \item corrected-matrix-element reconstruction; and
  \item validation reports separating matrix-element and observable gates.
\end{enumerate}

No production implementation should begin before the two-day literature sprint
freezes the first hypothesis, the exact comparison quantities, and the first two
material-mode contrasts.
